\section{Introduction}\label{sec:introduction}

MADNESS~\cite{harrison2016madness} is a multiresolution analysis (MRA) framework for density functional theory that represents functions as adaptive octree structures of polynomial scaling coefficients.
The tree refines locally where the function has fine-scale features and remains coarse elsewhere, achieving systematic accuracy without a fixed basis set.
Within a MADNESS DFT calculation, self-consistent field (SCF) iteration dominates wall time: each iteration rebuilds the Coulomb and exchange-correlation potentials, solves the bound-state Helmholtz equation for every orbital, and updates the density.
The standard initial density guess is the promolecular density $\rho_0$, a superposition of isolated spherical atomic densities.
Because $\rho_0$ contains no bonding, charge transfer, or exchange-correlation information, the SCF procedure typically requires 8--15 iterations to converge.
A more accurate starting density would reduce the iteration count and, proportionally, the total cost.

Machine learning models for electron density prediction have shown that neural networks can map nuclear geometry to real-space charge densities with useful accuracy~\cite{gong2025gedcrn}.
These approaches typically predict density on fixed Cartesian grids and project the result onto whatever basis the solver requires.
Working directly in the native MRA representation sidesteps grid-to-tree projection and preserves the adaptive structure that makes MADNESS efficient.
The octree further decomposes the prediction problem into per-node regression tasks with local spatial context available from same-level neighbor nodes, a natural inductive bias for an MLP architecture.

We present a systematic investigation of three ML strategies for accelerating SCF convergence in MADNESS: direct prediction of converged density coefficients, parent-augmented coefficient prediction, and tree structure (refinement) prediction.
All three use a FiLM-conditioned MLP (MRANet) with same-level halo neighbor features as input.
Despite strong per-node regression accuracy at fine tree levels, none of the strategies reduced SCF iteration count when the predicted density was used as an initial guess.
We trace this failure to an input signal bottleneck: at the coarse tree levels that determine SCF convergence behavior, the promolecular density $\rho_0$ and nuclear potential $V_\mathrm{nuc}$ lack the electron-electron interaction information needed to predict density corrections.
The model learns an accurate identity mapping at fine levels (where $\rho_0 \approx \rho$) but cannot extract the missing physics at coarse levels.
From this negative result we extract three lessons transferable to ML applied to multi-scale iterative solvers: (1) input signal quality at the scale that matters for the solver sets the model ceiling, regardless of aggregate accuracy; (2) fine-level accuracy does not imply solver acceleration; and (3) multi-task interference between level-dependent objectives can silently degrade training.
